\subsubsection{Representación rectilínea de las fases y del retorno}
\label{mono-recta:desarrollo}

La representación sobre el segmento \([0,10]\) conserva el origen
posicional de APP y ordena sus nueve marcas interiores. En este diagrama,
los enteros \(1,\ldots,9\) son representantes de fase. La posición
horizontal indica el orden de las coordenatizaciones; las etiquetas superiores
especifican operaciones regionales. La figura
\ref{mono-recta:figura} reúne así la recta de partida, la fijación del
eje de autoescala, la distribución especular de los dos canales y el
retorno con memoria (sección~\ref{mono-ampl:conexion}).

% Recuperación de fig_nonadic_route_mini.tex de la síntesis académica.
% Los rótulos designan operaciones regionales; la suma se efectúa en F_3^6.
\begin{figure}[H]
\centering
\begin{tikzpicture}[x=1.12cm,y=1cm,font=\small,>=Latex,
  guide/.style={line width=.45pt,black!55}]
 \draw[line width=.8pt] (0,0)--(10,0);
 \foreach \k in {0,...,10}{
  \draw[guide] (\k,-.09)--(\k,.09);
  \node[below=2.5mm] at (\k,0) {$\k$};
 }
 \foreach \k in {1,...,9}{\fill (\k,0) circle (1.7pt);}
 \node[anchor=west,font=\footnotesize] at (0,3.05)
   {Segmento generador $[0,10]$ y nueve representantes de fase};
 \draw[guide] (5,.12)--(5,1.7);
 \node[above,align=center,font=\footnotesize] at (5,1.7)
   {fijación del eje\\$u^\varphi$};
 \draw[decorate,decoration={brace,amplitude=4pt},guide]
   (6,.65)--(8,.65);
 \node[above=5pt,align=center,font=\footnotesize] at (7,.65)
   {agregado regional\\$u^e+u^\pi$};
 \draw[guide] (9,.12)--(9,1.7);
 \node[above,align=center,font=\footnotesize] at (9,1.7)
   {involución especular\\$\pi\leftrightarrow e$};
 \draw[->,line width=.7pt] (9,-.7)
   .. controls (9,-1.55) and (1,-1.55) .. (1,-.7);
 \node[fill=white,inner sep=3pt,font=\footnotesize] at (5,-1.23)
   {retorno $9\to1$;\quad $m\mapsto m+1$};
 \node[font=\footnotesize] at (5,-1.95)
   {$w_6\to w_{12}\to w_{18}\to w_{24}\to w_{30}\to R_{36}\to G_9$};
\end{tikzpicture}
\caption{Representación rectilínea de la conexión nonádica. La marca
\(5\) sitúa la primera saturación fuerte y el eje regional de
autoescala; la llave \(6\)--\(8\) identifica el agregado ternario
fijado en cada fibra; la fase \(9\) selecciona la orientación de
clausura. El retorno conserva la fase e incrementa la memoria. Las
abscisas numeran posiciones operativas, mientras las etiquetas
\(u^\pi,u^e,u^\varphi\) designan bloques de \(\mathbb F_3^6\).}
\label{mono-recta:figura}
\end{figure}


La fijación del eje se formula en la fibra de una frontera dada. Si
\(q=u^\pi+u^e+u^\varphi\) y
\(a=u^\pi+u^e-u^\varphi\), entonces
\[
 u^\varphi=2(q-a),\qquad s=u^\pi+u^e=q-u^\varphi
 \quad\text{en }\mathbb F_3^6.
\]
Los datos \((q,a)\) determinan el eje y el agregado; el reparto ordenado
de \(s\) conserva la libertad especular resuelta por la prolongación.
La quinta fase posee una solución local y una extensión, con dato
de frontera \(c=111111\) y firma residual \(\rho_W=000\). Éste es el
sentido de su primera saturación fuerte. El símbolo \(\varphi\)
situado sobre la marca \(5\) identifica el eje regional que interviene
en esa operación.

Las fases intermedias conservan esta descomposición dentro de cada fibra,
mientras sus datos de frontera evolucionan. El censo completo de
multiplicidades locales y extensiones de la rama distinguida es
\[
\begin{array}{c|rrrrrrrrr}
\text{fase}&1&2&3&4&5&6&7&8&9\\\hline
N_{\rm loc}&3&2&5&4&1&1&3&3&2\\
N_{\rm ext}&2&5&4&1&1&3&3&2&1
\end{array}
\]
La fase \(6\) combina unicidad local y tres prolongaciones; la fase
\(7\) resuelve una fibra especular triple; la fase \(8\) sincroniza
los tres censos regionales. La fase \(9\) conserva dos distribuciones
locales con el mismo eje y el mismo agregado, de las cuales una se
prolonga. Estos recuentos expresan funciones diferentes dentro de una
misma conexión, y justifican las marcas de la recta.

En particular, para las fases \(6,7,8\) las sumas regionales son
respectivamente \(012100\), \(101001\) y \(112000\). La conservación
representada por la llave es una invariancia respecto del reparto
especular en cada fibra. La evaluación real posterior aplica los
lectores de los cilindros compatibles a las historias completas; la
suma interna representada en la figura pertenece al dominio ternario.

Finalmente, el arco \(9\to1\) representa el paso de una vuelta a la
siguiente. Las fórmulas de la sección
\ref{mono-ampl:conexion} conservan el resto de fase e incrementan el
contador de vueltas. El diagrama rectilíneo y el diagrama cíclico de la
figura \ref{mono-ampl:fig-regional} expresan, por tanto, dos aspectos
complementarios de la misma holonomía: orden de las posiciones y
composición del retorno sobre el estado con memoria.
